\section{Variance from off-diagonal sampling}\label{sec:variance}

We now use the additive overlap identity to compare the mass of a cell
under a fiber law with its mass under the uniform law.  The same argument
will apply to the conditional law of sources arriving at a target.  It
is essential to sample distinct pairs throughout: the hierarchy has an
infinite range of singleton levels, so a small diagonal error cannot be
integrated indiscriminately.

\begin{lemma}\label{lem:variance}
Let $\mathcal F$ be a finite family of probability measures $\nu$, each
uniform on a subset $B\subset P$ of size $k\ge n^\gamma$.  Give the family
nonnegative weights, independent of places and grids, of total mass at
most one, and write $\E_{\mathcal F}$ for the weighted sum.  Suppose
$\E_{\mathcal F}\nu\le M\alpha$ pointwise, where $M$ is fixed.
For each transformed cell occurrence $C$, put
$m=|C|$, $a=\alpha(C)$, and $p=\nu(C)$.  Then
\begin{align}
\I\!\left(\E_{\mathcal F}\left[
  (p-a)^2\1_{\{m>n^{1-c}\}}
  +p^2\1_{\{m\le n^{1-c},\ p\ge2/k\}}
\right]\right)
&\le C_M\bigl(n^{c-\gamma}W_2+W_{2,\mathrm{sm}}\bigr)
\label{eq:variance-quant}\\
&=o(W+1).
\label{eq:variance}
\end{align}
The constant $C_M$ depends only on $M$.
\end{lemma}

The first term measures discrepancy from $\alpha$ on cells large enough
for the finite-population corrections to be negligible.  On a small
cell, the second term measures $\nu$-mass only when $B\cap C$ contains
at least two points: a one-point intersection contributes no
distinct-pair overlap within $B$.

\begin{proof}
The proof uses an exact signed identity for pair probabilities, followed
by separate estimates on large cells, small cells, and singletons.
The integer $k\ge n^\gamma>1$ is at least two.  Let $P_B$ be the uniform
law on distinct ordered pairs from $B$.  Define another law $R_B$ on
distinct pairs as follows.  First take $X$ uniformly from $B$.  If $k<n$,
then, conditionally on $X$, take $Y$ uniformly from $B\setminus\{X\}$
with probability $k/n$, and uniformly from $P\setminus B$ otherwise.
If $B=P$, define $R_B=P_B=\pi$ directly; the outside branch has zero
probability and need not be defined.

The first marginal of $R_B$ is $\nu$, and its second marginal is exactly
$\alpha$.  Indeed, a point $y\in B$ receives mass
\[
 (k-1)\frac1k\frac{k}{n}\frac1{k-1}=\frac1n,
\]
and a point $y\notin B$ receives mass
$k(1/k)((n-k)/n)/(n-k)=1/n$.  Thus the signed measure
$P_B-2R_B+\pi$ has total mass zero.  Its first and second marginals are
$\alpha-\nu$ and $\nu-\alpha$, respectively, so their sum is zero.
The additive overlap identity~\eqref{eq:additive} gives
\begin{equation}\label{eq:variance-cancel}
 \I\bigl(P_B(C^2)-2R_B(C^2)+\pi(C^2)\bigr)=0.
\end{equation}
This identity involves finite integrals.  For fixed $n$, any probability
law on distinct pairs is bounded by $n(n-1)\pi$ pointwise, so each of the
three nonnegative terms has a finite integral because $W_2<\infty$.
The family weights are independent of places and grids; hence we may
average~\eqref{eq:variance-cancel} over the family.

We calculate the three cell probabilities before making any estimates.
Writing $b=|B\cap C|$, so that $p=b/k$, gives
\begin{align*}
 P_B(C^2)
 &=\frac{b(b-1)}{k(k-1)}
   =p^2-\frac{p(1-p)}{k-1},\\
 R_B(C^2)
 &=p\left(\frac{k}{n}\frac{b-1}{k-1}+\frac{m-b}{n}\right)
   =ap-\frac{k}{n}\frac{p(1-p)}{k-1},\\
 \pi(C^2)
 &=\frac{m(m-1)}{n(n-1)}
   =a^2-\frac{a(1-a)}{n-1}.
\end{align*}
The middle identity also holds when $k=n$, by the direct definition of
$R_B$.  Therefore, with
$H_B(C)=P_B(C^2)-2R_B(C^2)+\pi(C^2)$,
\begin{equation}\label{eq:variance-finite}
 H_B(C)=(p-a)^2+\left(\frac{2k}{n}-1\right)
                 \frac{p(1-p)}{k-1}
                 -\frac{a(1-a)}{n-1}.
\end{equation}
Write $\overline H=\E_{\mathcal F}H_B$ and
$\omega=\E_{\mathcal F}1\le1$.  In the averaged formula, the last term
in~\eqref{eq:variance-finite} is multiplied by $\omega$.

First suppose $m>n^{1-c}$.  Since
$1/(k-1)\le2n^{-\gamma}$, $1/(n-1)\le2n^{-\gamma}$, and
$\E_{\mathcal F}p\le Ma$, Equation~\eqref{eq:variance-finite} implies
\[
 \overline H(C)\ge\E_{\mathcal F}(p-a)^2
                   -2(M+1)n^{-\gamma}a.
\]
For every $m\ge2$ we have
\begin{equation}\label{eq:variance-pair-comparison}
 \tfrac12a^2\le\pi(C^2)\le a^2.
\end{equation}
Here $a>n^{-c}$, so
$n^{-\gamma}a\le2n^{c-\gamma}\pi(C^2)$.  Hence
\begin{equation}\label{eq:variance-large}
 \overline H(C)\ge\E_{\mathcal F}(p-a)^2
       -4(M+1)n^{c-\gamma}\pi(C^2).
\end{equation}

Next suppose $2\le m\le n^{1-c}$.  If $p\ge2/k$, then $b\ge2$ and
\[
 P_B(C^2)=p^2\frac{k}{k-1}\frac{b-1}{b}\ge\tfrac12p^2.
\]
Since $P_B(C^2)$ is always nonnegative, this proves
$P_B(C^2)\ge\tfrac12p^2\1_{\{p\ge2/k\}}$ for every $B$.
Also $R_B(C^2)\le ap$, and therefore
\[
 \E_{\mathcal F}R_B(C^2)\le Ma^2\le2M\pi(C^2).
\]
Discarding the nonnegative term $\omega\pi(C^2)$ yields
\begin{equation}\label{eq:variance-small}
 \overline H(C)\ge\tfrac12\E_{\mathcal F}
                \bigl[p^2\1_{\{p\ge2/k\}}\bigr]-4M\pi(C^2).
\end{equation}

Finally, if $m=0$ or $m=1$, all three pair probabilities vanish exactly.
For a singleton, $p$ is either zero or $1/k$, so the small-cell term in
the conclusion vanishes too; a singleton cannot satisfy
$m>n^{1-c}$.  This exact cancellation treats the infinite singleton
tails.  We never integrate the separate correction terms
in~\eqref{eq:variance-finite} over those tails.

Integrate~\eqref{eq:variance-large} and~\eqref{eq:variance-small} on
their respective regions, use the exact zero on empty and singleton
cells, and apply~\eqref{eq:variance-cancel}.  The result is
\begin{align*}
 \I\!\left(\E_{\mathcal F}\left[
 (p-a)^2\1_{\{m>n^{1-c}\}}
 +\tfrac12p^2\1_{\{m\le n^{1-c},\ p\ge2/k\}}
 \right]\right)
 \le 4(M+1)n^{c-\gamma}W_2+4M W_{2,\mathrm{sm}}.
\end{align*}
The nonnegative left-hand terms are integrable by these pointwise
bounds and the integrability of $\overline H$ and the displayed errors.
Doubling the bound proves~\eqref{eq:variance-quant}, with
$C_M=8(M+1)$.

All the cell estimates hold for arbitrary subsets of $P$.
They therefore include giant complements and count repeated cell
occurrences with their original multiplicity.  In the special case
$B=P$, the signed expression and the large-cell variance are zero;
the small-cell term may equal $a^2$, which is absorbed by
$2\pi(C^2)$ as above.  Finally, $c<\gamma$, $W_2\le W$, and
Equation~\eqref{eq:small-overlap} give~\eqref{eq:variance}.
\end{proof}

We will use Lemma~\ref{lem:variance} for two families.  For the fiber
laws $\mu_e$, with weights $k_e/|E|$, the weights sum to one and
\[
 \E_e\mu_e=q_Y\le B_0\alpha.
\]
Their support sizes satisfy $k_e\ge n^{s-o(1)}\ge n^\gamma$ for all
sufficiently large $n$, since $\gamma<s$.

For the incoming laws $\rho_y$, retain the targets whose incoming
degree $d_y$ is at least $n^\gamma$, and use weights
$q_Y(y)=d_y/|E|$.  Their weighted mean is dominated by $q_X$:
\[
 \sum_{y:d_y\ge n^\gamma}q_Y(y)\rho_y\le q_X\le B_0\alpha.
\]
Each retained law is uniform on its $d_y$ incoming sources, and the
family weights have total mass at most one.  Both families are defined
by the same finite graph at every place, as required by the lemma.
The omitted targets have total mass
\begin{equation}\label{eq:variance-omitted}
 \sum_{y:d_y<n^\gamma}q_Y(y)
 \le\frac{n^{1+\gamma}}{|E|}
 =O(n^{\gamma-1})=o(1),
\end{equation}
because the edge density tends to $\theta>0$ and $\gamma<1$.
Thus the cell-mass estimates apply both to outgoing fibers and to
almost all incoming conditional laws.  These are the two distributions
that the remaining argument must compare with $\alpha$.
